\section{Centering vs.\ Direct Offset Landing}\label{sec:centering-analysis}

A critical design decision in the landing pipeline is whether to visually servo the drone above the target before commanding a landing offset, or to proceed directly from a GPS-estimated position.  The state machine includes a dedicated \textsc{Centering} state that flies the drone above the detected target, as well as a \textsc{Descending} state (currently bypassed) that would reduce altitude while maintaining visual lock (Section~\ref{sec:state-machine}). In the current implementation, the drone centres above the target at search altitude (\SI{35}{\metre}), enters VERIFY for operator confirmation and GPS averaging, then proceeds to an offset landing. This section analyses the centering-vs-direct-offset trade-off and justifies the current design choice.

% ─────────────────────────────────────────────────────────────────────
\subsection{The Centering Option}

In the full centering pipeline as originally designed, the drone transitions from \textsc{Search} to \textsc{Centering}, where it flies toward the GPS-estimated target position and attempts to place the detection bounding box within a pixel-tolerance of the image centre. The code includes a \textsc{Descending} state that would then reduce altitude from \SI{35}{\metre} to \SI{15}{\metre} in increments while maintaining visual lock. In the current implementation, the \textsc{Descending} state is bypassed: the drone transitions directly from \textsc{Centering} to \textsc{Verify} at search altitude (\SI{35}{\metre}). Only after the operator confirms the target does the system command a landing offset.

This pipeline improves localisation accuracy because lower-altitude frames have a smaller ground sampling distance (GSD), and centering eliminates the pixel-projection error from off-axis detections. However, it introduces three additional risk factors:

\begin{enumerate}
    \item \textbf{Low-altitude manoeuvring.} Descending from \SI{35}{\metre} to \SI{5}{\metre} while simultaneously correcting lateral position requires sustained GUIDED-mode velocity commands. Any interruption---communication dropout, momentary detection loss, or wind gust---can leave the drone at low altitude with stale position commands.
    \item \textbf{Visual servo fragility.} At altitudes below \SI{10}{\metre}, the target may exceed the camera field of view, motion blur increases with ground-relative speed, and the narrow depth of field of the IMX296 global shutter sensor reduces detection confidence. During simulation testing, three of twenty centering runs exhibited oscillatory lateral corrections at \SI{8}{\metre} altitude.
    \item \textbf{Increased flight time over the casualty.} The centering and descent phases add 30--60\,s of hover time directly above the target, which in a real SAR scenario places the casualty at greater risk from rotor downwash and potential mechanical failure.
\end{enumerate}

% ─────────────────────────────────────────────────────────────────────
\subsection{Why Direct Offset Landing Is Sufficient}

The offset landing approach bypasses centering: once the operator confirms the target, the drone flies to a GPS point displaced \SI{7.5}{\metre} from the estimated target position (perpendicular to the wind vector when available, otherwise due south) and executes a standard \texttt{MAV\_CMD\_NAV\_LAND}. The key insight is that the system requirement (R07) specifies landing \emph{within \SIrange{5}{10}{\metre}} of the casualty---not directly on top of it. A \SI{7.5}{\metre} offset places the nominal landing point at the centre of this annular band, and the question reduces to whether GPS-based localisation is accurate enough to keep the actual touchdown within the band.

% ─────────────────────────────────────────────────────────────────────
\subsection{Probability Analysis}\label{sec:prob-analysis}

The GPS-estimated target position was characterised in simulation with realistic noise parameters calibrated against DJI flight-video telemetry (Section~\ref{sec:localisation}). The circular error probable at the 50th percentile was measured as $\mathrm{CEP}_{50} = \SI{2.3}{\metre}$. Assuming the radial error follows a Rayleigh distribution, the scale parameter is

\begin{equation}\label{eq:sigma-rayleigh}
    \sigma = \frac{\mathrm{CEP}_{50}}{\sqrt{2\ln 2}} \approx \frac{2.3}{1.177} \approx \SI{1.95}{\metre}
\end{equation}

\noindent The actual landing point is offset by $d_0 = \SI{7.5}{\metre}$ from the estimated target. The distance from the \emph{true} target to the landing point is therefore a random variable whose distribution depends on the vector sum of the deterministic offset and the stochastic GPS error. In the worst case (all error components aligned radially), the landing distance from the true target is $r = d_0 + \epsilon$, where $\epsilon \sim \mathcal{N}(0, \sigma^2)$ projected onto the offset axis. The probability that the landing falls within the \SIrange{5}{10}{\metre} band is

\begin{equation}\label{eq:pband}
    P(5 < r < 10) \;=\; \Phi\!\left(\frac{10 - 7.5}{1.95}\right) - \Phi\!\left(\frac{5 - 7.5}{1.95}\right) \;=\; \Phi(1.28) - \Phi(-1.28) \;\approx\; 0.900 - 0.100 \;=\; 0.800
\end{equation}

\noindent where $\Phi$ is the standard normal CDF. This conservative uni-axial projection gives a lower bound of 80\%. In practice the error is two-dimensional: the lateral component (perpendicular to the offset direction) shifts the landing point around the annular band rather than outside it, and the multi-detection fusion (inverse-variance weighting over ${\sim}10$ observations) reduces $\sigma$ by $\sqrt{N_\text{eff}}$. With an effective sample size of $N_\text{eff} = 6$, the fused $\sigma_\text{fused} \approx 1.95 / \sqrt{6} \approx \SI{0.80}{\metre}$, giving

\begin{equation}\label{eq:pband-fused}
    P(5 < r < 10) \;=\; \Phi\!\left(\frac{2.5}{0.80}\right) - \Phi\!\left(\frac{-2.5}{0.80}\right) \;=\; \Phi(3.13) - \Phi(-3.13) \;\approx\; 0.999 - 0.001 \;>\; 0.99
\end{equation}

\noindent confirming that the fused GPS estimate with a \SI{7.5}{\metre} offset meets the R07 band requirement with over 99\% probability.

% ─────────────────────────────────────────────────────────────────────
\subsection{Simulation Comparison}

To validate the analysis, twenty simulated missions were executed for each approach using the SITL environment with GPS noise ($\sigma_\mathrm{GPS} = \SI{1.5}{\metre}$), camera shake, and \SI{3}{\metre\per\second} wind. Table~\ref{tab:centering-comparison} summarises the results.

\begin{table}[ht]
\centering
\caption{Comparison of landing approaches over 20 simulated missions each.}
\label{tab:centering-comparison}
\begin{tabular}{lcccc}
\toprule
\textbf{Approach} & \textbf{Mean Error} & \textbf{Max Error} & \textbf{In 5--10\,m Band} & \textbf{Notes} \\
\midrule
Direct offset (no centering) & \SI{3.1}{\metre} & \SI{4.8}{\metre} & 95\% (19/20) & 1 outlier at \SI{11.2}{\metre} \\
With centering + offset       & \SI{1.8}{\metre} & \SI{3.4}{\metre} & 100\% (20/20) & 3 runs showed low-alt instability \\
\bottomrule
\end{tabular}
\end{table}

\noindent The ``mean error'' column reports the absolute deviation of the actual landing distance from the intended \SI{7.5}{\metre} offset---i.e.\ how far the touchdown was from the centre of the acceptable band. The direct offset approach achieved 95\% compliance with R07, with a single outlier caused by a GPS multipath spike during the final descent. The centering approach improved mean accuracy to \SI{1.8}{\metre} but produced three runs where the visual servo exhibited lateral oscillation below \SI{10}{\metre} altitude, requiring the safety pilot to intervene (RC override to LOITER). In all three cases the drone recovered, but the behaviour would be unacceptable in an operational SAR deployment.

% ─────────────────────────────────────────────────────────────────────
\subsection{Design Decision}

Given that (a)~the direct offset approach meets R07 with $>$95\% probability under realistic noise, (b)~centering introduces low-altitude instability risks that are difficult to mitigate without extensive tuning, and (c)~the project timeline prioritises a reliable first demonstration over optimal accuracy, the direct offset landing was selected as the baseline configuration. The centering states remain implemented and tested in simulation; enabling them requires only a single configuration flag (\texttt{ENABLE\_CENTERING = True} in \texttt{config.py}), making this a low-risk upgrade path for future iterations.

% ─────────────────────────────────────────────────────────────────────
\subsection{Safety Considerations of Hovering Above the Casualty}\label{sec:hover-safety}

Visual centering requires the drone to hover directly above the detected casualty for 10--25\,s while the bounding box is driven toward the image centre. This raises a legitimate safety question: is it acceptable to position a UAV directly overhead an injured person? This subsection analyses the physical risks, quantifies them, compares mitigation strategies, and justifies the accepted tradeoff.

\subsubsection{Risk Identification}

Four hazards are associated with overhead hover:

\begin{enumerate}
    \item \textbf{Rotor downwash on the casualty.} Propeller wash could disturb an injured person, aggravate wounds, or displace lightweight medical dressings.
    \item \textbf{Mechanical failure leading to vertical descent onto the casualty.} A motor, ESC, or flight-controller failure at hover could cause the drone to fall directly onto the person below.
    \item \textbf{Psychological distress.} A conscious casualty may experience anxiety from a hovering, audible UAV directly overhead.
    \item \textbf{Debris entrainment.} Downwash may lift dust, leaves, or loose debris into the casualty's airway or wounds.
\end{enumerate}

\subsubsection{Quantitative Analysis}

\paragraph{Rotor downwash at ground level.}
The EDU-450 airframe has a rotor radius $R \approx \SI{0.225}{\metre}$ and an all-up mass of ${\sim}\SI{2}{\kilogram}$. In hover the induced velocity at the rotor disc is given by momentum theory:

\begin{equation}\label{eq:v-hover}
    v_h = \sqrt{\frac{T}{2\rho A}} = \sqrt{\frac{mg}{2\rho\pi R^2}} = \sqrt{\frac{2 \times 9.81}{2 \times 1.225 \times \pi \times 0.225^2}} \approx \SI{5.7}{\metre\per\second}
\end{equation}

\noindent where $T = mg$ is the hover thrust, $\rho = \SI{1.225}{\kilogram\per\metre\cubed}$ is air density, and $A$ is the total disc area of four rotors. The wake contracts and then expands as it descends; at a distance $h$ below the rotor plane the centreline velocity decays approximately as~\cite{johnson2013helicopter}:

\begin{equation}\label{eq:v-ground}
    v(h) \approx v_h \left(\frac{R}{h}\right)^2
\end{equation}

\noindent At the centering altitude of $h = \SI{35}{\metre}$:

\begin{equation}
    v_{\text{ground}} \approx 5.7 \times \left(\frac{0.225}{35}\right)^2 \approx 5.7 \times 4.1 \times 10^{-5} \approx \SI{0.0002}{\metre\per\second}
\end{equation}

\noindent This is four orders of magnitude below perceptible airflow (${\sim}\SI{0.5}{\metre\per\second}$) and six orders below the threshold for displacing lightweight objects (${\sim}\SI{3}{\metre\per\second}$). Even at the lower verification altitude of \SI{15}{\metre}:

\begin{equation}
    v_{15} \approx 5.7 \times \left(\frac{0.225}{15}\right)^2 \approx \SI{0.0013}{\metre\per\second}
\end{equation}

\noindent which remains entirely negligible. Rotor downwash is therefore a non-issue at any altitude above \SI{10}{\metre} for a \SI{2}{\kilogram} quadrotor.

\paragraph{Crash probability.}
Consider a total motor/ESC failure causing uncontrolled descent from \SI{35}{\metre}. Under free-fall with aerodynamic drag, the drone drifts laterally with the ambient wind during the ${\sim}\SI{2.7}{\second}$ fall time ($t = \sqrt{2h/g}$). In a \SI{5}{\metre\per\second} crosswind the lateral displacement is ${\sim}\SI{13}{\metre}$, giving a roughly circular impact zone of radius \SI{15}{\metre}. The probability that the drone strikes a casualty of cross-section ${\sim}\SI{1}{\metre\squared}$ within this $\pi \times 15^2 \approx \SI{707}{\metre\squared}$ area is:

\begin{equation}
    P_{\text{hit}} = \frac{A_{\text{casualty}}}{A_{\text{impact}}} \approx \frac{1}{707} \approx 0.14\%
\end{equation}

\noindent This is an upper bound (assumes the failure occurs at the worst moment and the casualty is at the centre). The base rate for catastrophic motor failure in a well-maintained quadrotor is ${\sim}10^{-4}$ per flight hour. For a 20\,s hover the conditional probability is $P_{\text{fail}} \approx 10^{-4} \times (20/3600) \approx 5.6 \times 10^{-7}$. The joint probability of failure \emph{and} casualty strike is therefore:

\begin{equation}
    P_{\text{fail \& hit}} \approx 5.6 \times 10^{-7} \times 1.4 \times 10^{-3} \approx 8 \times 10^{-10}
\end{equation}

\noindent which is negligibly small---comparable to the probability of being struck by lightning during the same interval.

\paragraph{Noise exposure.}
The EDU-450 produces approximately \SI{65}{\decibel} at \SI{1}{\metre}. At \SI{35}{\metre} the inverse-square law gives $65 - 20\log_{10}(35) \approx \SI{34}{\decibel}$, which is quieter than a whispered conversation (\SI{40}{\decibel}). Even at \SI{15}{\metre} the level is ${\sim}\SI{41}{\decibel}$---comparable to a quiet library.

\subsubsection{Comparison of Centering Strategies}

Table~\ref{tab:hover-safety} compares three centering strategies on safety, accuracy, and complexity.

\begin{table}[ht]
\centering
\caption{Comparison of centering strategies: safety risk, localisation accuracy, and implementation complexity.}
\label{tab:hover-safety}
\begin{tabular}{p{3.2cm}ccccp{3.8cm}}
\toprule
\textbf{Strategy} & \textbf{Downwash} & \textbf{Crash Risk} & \textbf{CEP$_{50}$} & \textbf{R07 Compliance} & \textbf{Notes} \\
\midrule
Direct overhead centering (\SI{35}{\metre})
    & Negligible & $8\times10^{-10}$ & \SI{1.5}{\metre} & $>$99\% & Best accuracy; standard SAR practice \\
\addlinespace
Lateral offset centering (\SI{10}{\metre} offset, \SI{35}{\metre} alt)
    & None & None & \SI{3.0}{\metre} & ${\sim}$95\% & Re-introduces GSD, focal-length, and yaw projection errors \\
\addlinespace
No centering (direct offset landing)
    & None & None & \SI{3.1}{\metre} & 95\% & Simplest; selected baseline (Section~\ref{sec:prob-analysis}) \\
\bottomrule
\end{tabular}
\end{table}

\noindent The lateral-offset strategy positions the drone \SI{10}{\metre} to the side of the target at the same altitude and infers the target's ground position from the off-axis pixel location using the GSD projection pipeline. While this eliminates overhead exposure entirely, it reintroduces all geometric error sources that centering is designed to eliminate: ground sampling distance uncertainty, focal length calibration error, and yaw-dependent east--north projection. The resulting CEP$_{50}$ of ${\sim}\SI{3}{\metre}$ is essentially identical to the no-centering baseline, negating the purpose of the manoeuvre while adding implementation complexity.

\subsubsection{Industry Practice}

Commercial SAR platforms routinely hover above search targets for visual assessment. The DJI Matrice 300 RTK, used by UK mountain rescue teams, hovers at 30--50\,m for thermal and optical inspection before directing ground teams~\cite{murphy2016disaster}. The CAA's CAP~722 guidance for small UAS operations does not prohibit overhead flight in emergency response contexts, and the JARUS SORA methodology classifies ground risk based on kinetic energy at impact rather than hover position. For a \SI{2}{\kilogram} drone the kinetic energy class is well within the ``low'' category regardless of position relative to persons on the ground.

\subsubsection{Accepted Tradeoff}

The analysis demonstrates that at \SI{35}{\metre} altitude the physical risks of overhead hover are negligible: downwash is unmeasurable, crash-and-hit probability is ${\sim}10^{-9}$, and noise is below conversational level. The accuracy benefit is substantial---CEP improves from \SI{3.1}{\metre} (no centering) to \SI{1.5}{\metre} (with centering)---which directly reduces the probability of violating R07 by landing too close to or too far from the casualty. In the SAR context, an inaccurate position estimate that causes the drone to land on top of the casualty is a \emph{more} dangerous outcome than a brief overhead hover at altitude. The centering approach is therefore retained as an available upgrade, and the safety case supports its deployment without additional mitigations beyond the existing \SI{35}{\metre} minimum centering altitude.
